\documentclass[11pt]{article}

\usepackage[margin=1in]{geometry}
\usepackage{amsmath,amssymb,amsthm,array,mathtools,microtype}
\usepackage[T1]{fontenc}
\usepackage[hidelinks]{hyperref}

\newtheorem{theorem}{Theorem}[section]
\newtheorem{proposition}[theorem]{Proposition}
\newtheorem{corollary}[theorem]{Corollary}
\newtheorem{lemma}[theorem]{Lemma}
\newcommand{\A}{\mathbb A}
\newcommand{\Spec}{\operatorname{Spec}}
\newcommand{\gdeg}{\operatorname{gdeg}}
\numberwithin{equation}{section}

\title{The Exact Nonproperness Locus of the Quadratic-Gauge Keller Maps}
\author{Roy van Rijn\\{\small Independent researcher}}
\date{17 November 2026}
\hypersetup{
  pdftitle={The Exact Nonproperness Locus of the Quadratic-Gauge Keller Maps},
  pdfauthor={Roy van Rijn},
  pdfsubject={Nonproperness and boundary sheets of polynomial Keller maps},
  pdfkeywords={Keller map, nonproperness locus, discriminant, Newton polygon,
    Jelonek locus}
}

\begin{document}
\maketitle

\begin{abstract}
For the all-degree quadratic-gauge Keller maps, we identify the reduced
nonproperness locus with the discriminant of the scalar inverse equation.
We compute the complete fiber table on the degenerate slice \(\Pi=0\) and
the exact global \(\Pi\)-adic order of the discriminant.  These geometric
results are independent of the prescribed finite-\'etale-fiber theorem.
\end{abstract}

\section{The quadratic-gauge input}

Let \(K\) be a characteristic-zero field and
\begin{equation}
 G(S)=g_1S+g_2S^2+\cdots+g_NS^N,\qquad g_1g_3g_N\neq0.
 \label{eq:seed}
\end{equation}
We use the quadratic gauge from the prescribed-fiber construction
\cite{Companion}, giving its formulas and reconstruction here.
In source coordinates put
\begin{equation}
 t=1+xy,\qquad q=t^2z+\frac{g_1}{g_3}y^2(1+3t).
 \label{eq:source-tq}
\end{equation}
Define the polynomial map \(F_G=(\Pi,B,C):\A^3_K\to\A^3_K\) by
\begin{align}
 \Pi={}&tq,\notag\\
 B={}&y+3\frac{g_3}{g_1}xq+2\frac{g_2}{g_1}tq
       +\sum_{k=4}^N k\frac{g_k}{g_1}t^2x^{k-2}q^k,\notag\\
 C={}&x(5-3t)-\frac{g_3}{g_1}x^3z
       -\sum_{k=4}^N(k-2)\frac{g_k}{g_1}(xq)^k.
 \label{eq:source-map}
\end{align}
Its scalar inverse equation is
\[
 E_{\Pi,B,C}(S)=G_\Pi(S)-\frac{g_1}{2}(BS^2+C),
\]
where
\[
 G_\Pi(S)=g_1S+\Pi(g_2S^2+g_3S^3)
          +\sum_{k=4}^Ng_k\Pi^kS^k.
\]
It satisfies
\[
 \det DF_G=-2,\qquad \gdeg(F_G)=N.
\]
For \(\pi\neq0\), its literal fiber is
\begin{equation}
 F_G^{-1}(\pi,b,c)\simeq
 \Spec\left((K[S]/(E_{\pi,b,c}))[1/\overline{E_{\pi,b,c}'}]\right).
 \label{eq:localized-input}
\end{equation}
For clarity, the compact calculation verifying these assertions is as
follows.  On \(t\neq0\), set
\[
 S=x/t,\quad Q=y+xq,\quad D=1-SQ+\Pi S^2=t^{-1},\quad
 \beta(\Pi,S)=\frac{G_\Pi'(S)/g_1-1-\Pi S^2}{S}.
\]
The last quotient is a polynomial.  Substituting \eqref{eq:source-tq} into
\eqref{eq:source-map}, coefficient by coefficient, gives
\[
 B=Q+\beta(\Pi,S),\qquad C=2G_\Pi(S)/g_1-BS^2,\qquad
 D=E'_{\Pi,B,C}(S)/g_1.
\]
Direct differentiation gives
\[
 \det\frac{\partial(\Pi,S,Q)}{\partial(x,y,z)}=t=D^{-1},\qquad
 \det\frac{\partial(\Pi,B,C)}{\partial(\Pi,S,Q)}=-2D.
\]
Consequently the polynomial Jacobian determinant is \(-2\).
For \(\pi\neq0\), every fiber point has \(t,q\) invertible because
\(tq=\pi\).  Conversely, a root \(s\) with derivative unit reconstructs by
\[
 \begin{gathered}
 d=E'_{\pi,b,c}(s)/g_1,\quad Q=b-\beta(\pi,s),\quad
 t=d^{-1},\quad x=s/d,\quad y=Q-\pi s,\quad q=\pi d,\\
 z=d^2\left(q-\frac{g_1}{g_3}y^2(1+3/d)\right).
 \end{gathered}
\]
The identity \(d=1-sQ+\pi s^2\) verifies \(1+xy=t\); the other identities
then verify \(F_G(x,y,z)=(\pi,b,c)\).  These operations are inverse over
any commutative test algebra, proving \eqref{eq:localized-input}.
They also give the rational function-field equality
\(K(x,y,z)=K(\Pi,B,S)\), with \(\Pi,B,S\) algebraically independent and
\(C=2G_\Pi(S)/g_1-BS^2\) of degree \(N\) in \(S\).  This polynomial
has degree \(N\) over \(K(\Pi,B)\), so the induced extension
\(K(\Pi,B,S)/K(\Pi,B,C)\) has degree \(N\), proving the generic degree.

The boundary formulas used below follow from the same polynomial map.
On \(q=0\), one has \(\Pi=0\) and \(B=y\).  This divisor has
\(t\neq0\), since at \(t=0\) the equation for \(q\) is
\(q=(g_1/g_3)y^2\neq0\).  Thus \(S=x/t\) is defined there, and the compact
identities give
\[
 C=2S-BS^2,\qquad 1-BS=t^{-1}.
\]
On \(t=0\), one has \(xy=-1\),
\(q=(g_1/g_3)y^2\), and \(B=-2y=2/x\).  In the expression for \(C\),
the coefficient of \(z\) is \(-(g_3/g_1)x^3\neq0\).  These identities
suffice to check the complete boundary table without the companion paper.

For a squarefree \(P\in K[T]\) of degree \(N\geq3\) and
\(a\in K\) with \(P'(a)P'''(a)\neq0\), the prescribed-fiber seed is
\[
 G(S)=P(a+S)-P(a).
\]
For the displayed map \(F_G\), the distinguished target is
\begin{equation}
 y_{P,a}=\left(1,0,-\frac{2P(a)}{P'(a)}\right),
 \label{eq:prescribed-target}
\end{equation}
and its inverse equation is \(P(a+S)\).  If determinant one is required,
choose the output automorphism \(L(\Pi,B,C)=(\Pi,-B/2,C)\),
so that \(\det D(L\circ F_G)=1\).  Since the target has \(B=0\),
its coordinates are unchanged by \(L\).  The discriminant and boundary
formulas below use the displayed, unnormalized coordinates of \(F_G\);
the corresponding loci for \(L\circ F_G\) are their images under \(L\).

\section{The exact nonproperness locus}
\label{sec:nonproperness}

The prescribed-fiber realization theorem does not require the calculation
below.  This companion note gives an intrinsic geometric description of
fullness for the quadratic gauge.

Work geometrically over
\(\overline K\).  For a polynomial map \(F:\A^n_{\overline K}\to
\A^n_{\overline K}\), let \(S_F\) be the set of target points at which \(F\)
is not finite over any Zariski neighborhood.  Since source and target are
affine, this is the same as failure of properness over every neighborhood.
Equivalently, compactify the graph in
\(\mathbb P^n_{\overline K}\times\A^n_{\overline K}\); then \(S_F\) is the
image of its boundary.  We give \(S_F\) its reduced closed-subscheme
structure.  Over \(\mathbb C\) this agrees with the usual analytic Jelonek
set, defined by sequences escaping to infinity.

For a polynomial \(h\), write \(V_{\mathrm{red}}(h)\) for the reduced closed
subscheme \(V(\sqrt{(h)})\).  Put
\begin{equation}
 \Delta_G(\Pi,B,C)
 =\operatorname{Disc}_S\bigl(E_{\Pi,B,C}(S)\bigr).
 \label{eq:gauge-discriminant}
\end{equation}
Here the discriminant is the universal degree-\(N\) polynomial expression in
the coefficients of \(E_{\Pi,B,C}\).  We do not replace it on a special
slice by the discriminant of the lower-degree specialization.

\begin{lemma}[Base change of the graph-boundary locus]
\label{lem:nonfinite-base-change}
Let \(k\subset k'\) be a field extension and let
\(F:\A^n_k\to\A^m_k\) be a polynomial map.  Let \(\overline\Gamma_F\) be
the scheme-theoretic closure of its graph in
\(\mathbb P^n_k\times\A^m_k\), put
\[
 Z_F=\overline\Gamma_F\cap
 \bigl((\mathbb P^n_k\setminus\A^n_k)\times\A^m_k\bigr),
\]
and give the image of the proper morphism \(Z_F\to\A^m_k\) its reduced
closed-subscheme structure \(S_F\).  Then
\[
 Z_{F_{k'}}=Z_F\times_k k',
 \qquad
 \lvert S_{F_{k'}}\rvert=\lvert (S_F)_{k'}\rvert .
\]
In particular, whether a geometric target point lies in the graph-boundary
locus is invariant under extension of algebraically closed fields.
\end{lemma}

\begin{proof}
Field extension is faithfully flat.  On each standard affine chart of
\(\mathbb P^n_k\times\A^m_k\), the ideal of the graph closure is the
saturation \(I:x_0^\infty\) of the graph ideal by the equation \(x_0\) of
the hyperplane at infinity.  Since the coordinate rings are noetherian, this
saturation equals \(I:x_0^r\) for some \(r\); flatness gives
\[
 (I:x_0^r)\otimes_k k'=(I\otimes_k k'):x_0^r.
\]
Thus scheme-theoretic graph closure commutes with the field extension.
Intersecting with \(x_0=0\) also commutes with base change, which proves the
identity for \(Z_F\).

For the image statement, let \(y'\) be a point of \(\A^m_{k'}\), with image
\(y\) in \(\A^m_k\).  The fiber \((Z_F)_{y'}\) is obtained from
\((Z_F)_y\) by a field extension.  A finite-type scheme over a field is
empty if and only if it becomes empty after a field extension.  Hence
\((Z_F)_{y'}\) is nonempty exactly when \((Z_F)_y\) is nonempty.  Since a
proper morphism has closed image, these fiberwise equivalences give the
asserted equality of underlying closed sets.
\end{proof}

\begin{lemma}[Fiber cardinality and nonproperness in characteristic zero]
\label{lem:properness-cardinality}
Let \(k\) be an algebraically closed field of characteristic zero and let
\(F:\A^n_k\to\A^n_k\) be a dominant \'etale polynomial map of geometric
degree \(d\).  Then
\begin{equation}
 y\notin S_F\quad\Longleftrightarrow\quad
 \#F^{-1}(y)=d.
 \label{eq:properness-cardinality}
\end{equation}
\end{lemma}

\begin{proof}
Jelonek proves this multiplicity criterion for complex polynomial maps
\cite[Proposition~6, p.~261]{Jelonek1993}.  We record why that stated complex result
is sufficient here.  Since \(F\) is \'etale, its Jacobian determinant is a
nonzero unit of \(k[x_1,\ldots,x_n]\), hence a nonzero constant.  Descend the
coefficients of \(F\), this constant Jacobian identity, and the coordinates
of \(y\) to a finitely generated field \(k_0/\mathbb Q\).  The resulting map
\(F_0\) is \'etale and therefore dominant; its geometric degree is still
\(d\).  We also obtain \(y_0\in\A^n(k_0)\) whose base changes are \(y\).

Choose an embedding \(k_0\hookrightarrow\mathbb C\) and an algebraic closure
\(\Omega\) of \(k_0\).  The two embeddings of \(k_0\), into \(k\) and into
\(\mathbb C\), extend to embeddings
\[
 \Omega\hookrightarrow k,\qquad \Omega\hookrightarrow\mathbb C.
\]
Lemma~\ref{lem:nonfinite-base-change} shows that membership of \(y_0\) in
the graph-boundary locus has the same truth value after each of these scalar
extensions.  The fiber of an \'etale morphism is finite and reduced; its
number of geometric points is unchanged by an extension of algebraically
closed fields.  The degree of the generic finite separable extension is
likewise invariant under scalar extension.  We may therefore test both
sides of \eqref{eq:properness-cardinality} after base change to
\(\mathbb C\).  There Proposition~6 of \cite{Jelonek1993}, with every local
multiplicity equal to one, gives the equivalence.
\end{proof}

\begin{lemma}[Exact \(\Pi\)-adic discriminant order]
\label{lem:exact-discriminant-order}
Assume \(N\geq4\), put \(r=N-3\), and let
\[
 R=\overline K[B,C],\qquad L=\operatorname{Frac}(R).
\]
For the degree-\(N\) polynomial \(E=E_{\Pi,B,C}(S)\), viewed in
\(L((\Pi))[S]\), extend \(\operatorname{ord}_\Pi(\Pi)=1\) to an algebraic
closure of \(L((\Pi))\).  The roots split into three Newton blocks:
\begin{equation}
\begin{array}{c|c|c}
\text{number of roots}&\operatorname{ord}_\Pi(s)&
 \text{residual polynomial}\\ \hline
2&0&BZ^2-2Z+C\\
1&-1&Z-\dfrac{g_1B}{2g_3}\\[4pt]
r&-\dfrac{r+2}{r}&g_3+g_NZ^r.
\end{array}
\label{eq:newton-blocks}
\end{equation}
All three residual polynomials are separable over \(L\).  Consequently
\[
 \operatorname{ord}_\Pi\operatorname{Disc}_S(E)
 =N^2-3N-2,
\]
and, in \(R[\Pi]\),
\[
 \operatorname{Disc}_S(E)=\Pi^{N^2-3N-2}D_G,\qquad
 D_G(0,B,C)=
 \pm\frac{r^r}{4}g_1^4g_3^{r+1}g_N^{r-1}B^2(1-BC).
\]
\end{lemma}

\begin{proof}
The coefficient valuations give the lower Newton polygon
\[
 (0,0)\longrightarrow(2,0)\longrightarrow(3,1)
 \longrightarrow(N,N).
\]
Its slopes are \(0\), \(1\), and \((r+2)/r\), with horizontal lengths
\(2\), \(1\), and \(r\).  Rescaling \(S\) along each edge gives exactly the
three residual polynomials in \eqref{eq:newton-blocks}.  The first has
discriminant \(4(1-BC)\neq0\) in \(L\), the second is linear, and the third
is separable because \(\operatorname{char}L=0\) and \(g_3g_N\neq0\).

Write the roots in the three blocks as \(A_1,A_2\), \(T\), and
\(V_1,\ldots,V_r\).  Roots in different blocks have different valuations,
so the valuation of their difference is the smaller of the two
valuations.  Within the first and third blocks, distinct roots of the
residual polynomial have distinct initial terms.  Thus
\[
 \operatorname{ord}_\Pi(A_1-A_2)=0,\qquad
 \operatorname{ord}_\Pi(V_i-V_j)=-\frac{r+2}{r}\quad(i\neq j).
\]
This is the precise use of residual separability.  It proves equality in
each entry of the following ledger:
\[
\begin{array}{c|c|c|c}
\text{pair type}&\text{number of pairs}&
\operatorname{ord}_\Pi(s_i-s_j)&
\text{contribution to }\operatorname{ord}_\Pi\prod_{i<j}(s_i-s_j)^2\\ \hline
A\!-\!A&1&0&0\\
A\!-\!T&2&-1&-4\\
A\!-\!V&2r&-(r+2)/r&-4(r+2)\\
T\!-\!V&r&-(r+2)/r&-2(r+2)\\
V\!-\!V&r(r-1)/2&-(r+2)/r&-(r-1)(r+2).
\end{array}
\]
The leading coefficient of \(E\) is \(g_N\Pi^N\), so its contribution to
\[
 \operatorname{Disc}(E)
 =\operatorname{lc}(E)^{2N-2}\prod_{i<j}(s_i-s_j)^2
\]
is \(2N(N-1)\).  Adding the last column gives
\[
 2N(N-1)-4-4(r+2)-2(r+2)-(r-1)(r+2)
 =N^2-3N-2.
\]

It remains only to identify the first nonzero coefficient.  Let
\(\alpha_1,\alpha_2\) be the roots of \(BZ^2-2Z+C\), put
\(b_0=g_1B/(2g_3)\), and let \(v_1,\ldots,v_r\) be the roots of
\(g_3+g_NZ^r\).  Taking the initial term in every factor of the root-product
formula gives
\begin{align*}
 &\pm g_N^{2N-2}
 (\alpha_1-\alpha_2)^2 b_0^4
 \left(\prod_jv_j\right)^6
 \prod_{i<j}(v_i-v_j)^2\\
 &=\pm\frac{r^r}{4}
 g_1^4g_3^{r+1}g_N^{r-1}B^2(1-BC),
\end{align*}
because
\[
 (\alpha_1-\alpha_2)^2=\frac{4(1-BC)}{B^2},\qquad
 \left(\prod_jv_j\right)^6=\left(\frac{g_3}{g_N}\right)^6,
\qquad
 \prod_{i<j}(v_i-v_j)^2
 =\pm r^r\left(\frac{g_3}{g_N}\right)^{r-1}.
\]
This coefficient is nonzero in \(R\).  Since \(\Pi\)-adic order in the
domain \(R[\Pi]\) is unchanged after passage to \(L[\Pi]\), the asserted
factorization and specialized coefficient follow.
\end{proof}

\begin{theorem}[Exact reduced quadratic-gauge nonproperness]
\label{thm:exact-nonproperness}
For every seed \eqref{eq:seed},
\begin{equation}
 \boxed{S_{F_G}=V_{\mathrm{red}}(\Delta_G).}
 \label{eq:jelonek-discriminant}
\end{equation}
Equivalently, their sets of geometric points satisfy
\(\lvert S_{F_G}\rvert=V_{\overline K}(\Delta_G)\).
More precisely, the following statements hold.
\begin{enumerate}
\item If \(\pi\neq0\) and the multiplicity partition of
      \(E_{\pi,b,c}\) is \((m_1,\ldots,m_s)\), then
      \begin{equation}
       \#F_G^{-1}(\pi,b,c)=\#\{i:m_i=1\},
       \qquad
       N-\#F_G^{-1}(\pi,b,c)=\sum_{m_i\geq2}m_i.
       \label{eq:partition-sheet-loss}
      \end{equation}
\item On \(\Pi=0\), the complete geometric fiber table is
      \begin{equation}
      \begin{array}{c|c|c}
       \text{target stratum}&\#F_G^{-1}(0,B,C)&\text{lost sheets}\\ \hline
       B\neq0,\ BC\neq1&3&N-3\\
       B\neq0,\ BC=1&1&N-1\\
       B=0&1&N-1.
      \end{array}
      \label{eq:pi-zero-fiber-table}
      \end{equation}
\item For \(N=3\), the raw resultant has an extraneous leading-coefficient
      factor:
      \[
       \operatorname{Res}_S(E,E')=\text{\rm unit}\cdot\Pi\Delta_G,
       \qquad
       \Delta_G(0,B,C)=\frac{g_1^4}{4}B^2(1-BC).
      \]
      Hence the entire plane \(\Pi=0\) is not a nonproperness component in
      degree three.
\item For \(N\geq4\), put
      \[
       e_N=N^2-3N-2.
      \]
      Then
      \begin{equation}
       \Delta_G=\Pi^{e_N}D_G,
       \qquad
       D_G(0,B,C)=\kappa_G B^2(1-BC),
       \qquad \kappa_G\in K^*.
       \label{eq:saturated-discriminant}
      \end{equation}
      In particular,
      \[
       S_{F_G}=V_{\mathrm{red}}(\Pi D_G),
      \]
      whose underlying set is \(V(\Pi)\cup V(D_G)\), and the resultant and
      discriminant define the same reduced closed subscheme.
\end{enumerate}
\end{theorem}

\begin{proof}
Because \(F_G\) is \'etale, every geometric fiber is finite and reduced.
Its generic cardinality is \(N\), so
Lemma~\ref{lem:properness-cardinality} applies with \(d=N\).

Suppose first that \(\pi\neq0\).  The localized reconstruction
\eqref{eq:localized-input} identifies the fiber with
\[
 \operatorname{Spec}
 \bigl((\overline K[S]/(E_{\pi,b,c}))[1/\overline{E_{\pi,b,c}'}]\bigr).
\]
Thus a distinct root contributes one affine point exactly when it is simple.
This proves \eqref{eq:partition-sheet-loss}; it also shows that the fiber has
cardinality \(N\) exactly when \(\Delta_G(\pi,b,c)\neq0\).

It remains to analyze \(\Pi=0\), where the inverse chart used above is not
available.  In the source, \(\Pi=tq\).  On \(q=0\), one has
\[
 B=y,\qquad C=2S-BS^2,\qquad S=\frac{x}{t}.
\]
The residual inverse equation is therefore
\begin{equation}
 BS^2-2S+C=0.
 \label{eq:q-zero-residual}
\end{equation}
A root reconstructs on \(q=0\) precisely when \(1-BS\neq0\).  Hence this
divisor contributes two points when \(B\neq0\) and \(BC\neq1\), no point when
\(B\neq0\) and \(BC=1\), and one point when \(B=0\).

On \(t=0\), the relation \(1+xy=0\) gives \(y=-1/x\), while
\[
 q=\frac{g_1}{g_3}y^2,\qquad B=-2y=\frac2x.
\]
Thus \(t=0\) contributes one point when \(B\neq0\) and none when \(B=0\).
Once \(x\) and \(y\) are fixed, the equation for \(C\) is linear in \(z\)
with nonzero coefficient, so the point is unique.  The divisors \(t=0\) and
\(q=0\) do not meet.  This proves \eqref{eq:pi-zero-fiber-table}.

For \(N=3\), specializing the degree-three discriminant at \(\Pi=0\) gives
\[
 \Delta_G(0,B,C)=\frac{g_1^4}{4}B^2(1-BC).
\]
Its zero set is exactly the two short-fiber strata in
\eqref{eq:pi-zero-fiber-table}.  Since the leading coefficient is
\(g_3\Pi\), the resultant has one additional factor \(\Pi\).

Now let \(N\geq4\) and put \(r=N-3\).
Lemma~\ref{lem:exact-discriminant-order} gives the global polynomial
identity
\[
 \Delta_G=\Pi^{e_N}D_G,\qquad
 D_G(0,B,C)=\kappa_GB^2(1-BC),
\quad
 \kappa_G=\pm\frac{r^r}{4}
 g_1^4g_3^{r+1}g_N^{r-1}\neq0.
\]
In particular, \(e_N\) is the exact global \(\Pi\)-adic order.  After specializing
\((B,C)\), the pointwise order is \(e_N\) when
\(B(1-BC)\neq0\) and may jump on the two displayed special strata.

No general assertion about discriminants under a drop of leading coefficient
is being used.  For \(N=3\), the explicit specialized formula cuts out
exactly the two short-fiber strata.  For \(N\geq4\), \(\Delta_G\) vanishes on
the whole slice \(\Pi=0\), and every fiber on that slice has fewer than \(N\)
points by \eqref{eq:pi-zero-fiber-table}.  On \(\Pi\neq0\), the ordinary
degree-\(N\) discriminant detects exactly the repeated-root, hence
short-fiber, locus.
The fiber table and \eqref{eq:properness-cardinality} now show that the zero
set of \(\Delta_G\) is exactly the underlying nonproperness set, and taking
radicals gives \eqref{eq:jelonek-discriminant}.
\end{proof}

\begin{corollary}[Intrinsic fullness of the distinguished fiber]
\label{cor:intrinsic-fullness}
For the seed \(G(S)=P(a+S)-P(a)\) associated with a squarefree
polynomial presentation, the distinguished target
\eqref{eq:prescribed-target} lies in the
maximal open subset on which \(F_G\) is finite \'etale of degree
\(N\).
\end{corollary}

\begin{proof}
At \(y_{P,a}\), the inverse polynomial is \(P(a+S)\).  Hence
\[
 \Delta_G(y_{P,a})=\operatorname{Disc}(P(a+S))
 =\operatorname{Disc}(P)\neq0.
\]
Theorem~\ref{thm:exact-nonproperness} places the target outside the
nonproperness hypersurface.  Proper and \'etale is finite \'etale, and the
degree is \(N\).
\end{proof}


\section*{Verification}

The graph-boundary base-change argument, Jelonek descent, complete
\(\Pi=0\) fiber table, and Newton-polygon discriminant calculation are
ordinary mathematical proofs and are not part of the Lean formalization of
the prescribed-fiber theorem.  The accompanying source bundle includes a symbolic checker for selected
seeds, boundary identities and discriminant orders.  Its README states the
command and finite input range.  These checks supplement the all-degree
proof and do not verify it formally.  The versioned companion source
\cite{Companion} provides the broader prescribed-fiber context.

\paragraph{AI assistance and review.}
ChatGPT (OpenAI) assisted with mathematical arguments, literature checks,
exposition, proof review and verification code.

\begin{thebibliography}{9}
\footnotesize

\bibitem{Companion}
R.~van Rijn, \emph{Over Characteristic Zero, Every Finite \'Etale Algebra
of Rank at Least Three Is a Full Keller Fiber},
repository edition \texttt{c32b9b4}, 17 November 2026,
\href{https://github.com/royvanrijn/jacobian-research/tree/c32b9b4d273cf7988e85184a0d5eae9c488727fd/unpublished/common-arithmetic-fibers}{versioned manuscript and source}.

\bibitem{Jelonek1993}
Z.~Jelonek,
\emph{The set of points at which a polynomial map is not proper},
Ann. Polon. Math. \textbf{58} (1993), 259--266,
\href{https://doi.org/10.4064/ap-58-3-259-266}
{doi:10.4064/ap-58-3-259-266}.

\end{thebibliography}

\end{document}
